\documentclass[11pt]{article}
\usepackage[margin=1in]{geometry}
\usepackage{amsmath,amssymb,amsthm}
\usepackage[hidelinks]{hyperref}
\setlength{\emergencystretch}{1em}
\newtheorem{proposition}{Proposition}
\newcommand{\R}{\mathbb{R}}
\newcommand{\HH}{\mathcal{H}}
\title{Disproof of conjecture 00000002531\\\large The stated density equals two on the real line}
\author{C0ldSmi1e}
\date{October 5, 2026}
\begin{document}
\maketitle

\paragraph{Claim and conventions.}
Both language versions define the density by
\[
 \theta_E(x)=\lim_{r\downarrow0}\frac{\HH^k(E\cap B(x,r))}{r^k}
\]
and assert that every $\HH^k$-measurable set has, outside an $\HH^k$-null
set, a density limit in $\{0,1\}$. They also assert dimension and sharpness
bounds for the set of points where the limit fails to exist. The denominator
above is exactly the source's $r^k$.

The source leaves the ambient space and normalization implicit. We use the
ordinary Euclidean interpretation, positive integer dimension, open balls
centered at $x$, and all positive real radii approaching zero. Standard
normalized Hausdorff measure is $\omega_k/2^k$ times the diameter-cover
measure, where $\omega_k$ is the volume of the Euclidean unit ball
\cite{Simon}. We take $k=n=1$. Then $\omega_1=2$, so the coefficient is
one, and $\HH^1$ on $\R$ is ordinary length. These conventions are explicit
in the Lean definitions. No unit-ball factor is inserted into the source's
density denominator.

\begin{proposition}
Let $E=\R$ with its usual distance $d(x,y)=|x-y|$. With the stated density
formula and standard $\HH^1$, the density exists and equals $2$ at every
$x\in\R$. No $\HH^1$-null exceptional set makes this density binary.
\end{proposition}
\begin{proof}
The set $E$ is Borel and hence $\HH^1$-measurable. For every $x\in\R$
and every $r>0$,
\[
 E\cap B(x,r)=(x-r,x+r),\qquad
 \HH^1(E\cap B(x,r))=2r,\qquad
 \frac{\HH^1(E\cap B(x,r))}{r^1}=2.
\]
Thus the full positive-radius limit is $2$, which is neither $0$ nor $1$.
If a null set $N$ made the binary conclusion true for every $x\in E\setminus N$,
then $E\setminus N$ would be empty. This forces $N=\R$, contradicting
$\HH^1(N)=0$, since $\HH^1(\R)=\infty$.
\end{proof}

\paragraph{Consequence and scope.}
The first necessary assertion is false, so the conjunction in the conjecture
is false. This example has a limit everywhere; it does not separately refute
the later dimension or sharpness assertions. The source imposes no finite
total measure hypothesis. Although $\HH^1(E)=\infty$, every ball numerator
in the displayed calculation is finite. The disproof applies whether
``almost everywhere'' is read on $E$ or on the ambient space, since here
$E=\R$. It depends on the stated normalization and denominator; replacing
the denominator by $\omega_k r^k$ would be a different assertion.

\newpage
\paragraph{Correspondence with Lean.}
The project pins Lean 4.19.0 and Mathlib v4.19.0. All mathematics is in
\texttt{Conjecture2531.lean}, under the namespace \texttt{Conjecture2531}.
The genuine library measure \texttt{Measure.hausdorffMeasure 1} is multiplied
by the explicitly defined coefficient $\omega_1/2^1$.
Theorems prove this coefficient is one and the resulting measure equals
\texttt{volume}, using Mathlib's \texttt{hausdorffMeasure\_real}. The genuine
real metric, ball measure, finite ball measures, and infinite total measure
are also checked.

The function \texttt{densityQuotient} divides the ball-intersection measure
by the extended nonnegative real value of $r^1$. Its extended codomain
preserves infinity. For $r>0$, the denominator is positive and finite.

\texttt{HasBinaryDensity} asserts convergence to a value equal to zero or one
along \texttt{nhdsWithin 0 (Set.Ioi 0)}. The radius filter is proved
nontrivial; the theorem uses uniqueness of limits, with the usual topology
on the extended nonnegative reals.

\texttt{NullExceptionalBinaryDensity} quantifies an arbitrary null set $N$
for each $E$ and requires binary density at every point of $E\setminus N$.
For Borel $E$, a theorem identifies this with the restricted-measure
almost-everywhere formulation, without assuming the density predicate is
measurable. A separate theorem shows that ambient almost-everywhere binary
density implies this same necessary assertion.

The definition \texttt{HausdorffMeasurable} uses the actual Carath\'eodory
predicate of the Hausdorff outer measure. Borel measurability implies it.
The formal propositions \texttt{BorelDensityClause} and
\texttt{CaratheodoryDensityClause} are explicitly necessary specializations
of the first source assertion to $k=n=1$. They are not definitions of the
entire compound conjecture. The main results are:
\begin{itemize}
\item \texttt{counterexample}: $\R$ is measurable in both senses, its density
converges to two everywhere, and no null exception restores binary density;
\item \texttt{not\_borelDensityClause} and
\texttt{not\_caratheodoryDensityClause}: negations of the necessary clauses;
\item \texttt{not\_statement\_implying\_borelDensityClause}: a logical bridge
from any explicitly supplied necessary-clause implication to its negation;
\item \texttt{not\_borelDensityClause\_and}: further conjunctive assertions
cannot rescue the false clause; those assertions remain uninterpreted.
\end{itemize}

\paragraph{Verification.}
A fresh independent build and source replay with warnings treated as errors
passed. The audit covers 31 authored declarations (nine definitions and
22 theorems) and six generated declarations, including their logical
dependencies. There are no admitted proofs, unsafe or partial dependencies,
custom axioms, or \texttt{native\_decide} calls. All axioms belong to the
standard set: \texttt{propext}, \texttt{Classical.choice}, \texttt{Quot.sound}.
The auxiliary inspection replay passed. No separate numerical experiment
is required. Detailed records and reproduction instructions are included;
local verification does not constitute maintainer acceptance.

\begin{thebibliography}{9}
\bibitem{Simon} Leon Simon, \emph{Introduction to Geometric Measure Theory},
Chapter 1, \S2 (Hausdorff measure normalization),
\href{https://math.stanford.edu/~lms/ntu-gmt-text.pdf}{author-hosted text}.
The source conjecture's density denominator is retained independently.
\bibitem{Mathlib} Mathlib v4.19.0,
\texttt{Mathlib.MeasureTheory.Measure.Hausdorff} and the real Lebesgue
measure library; exact revision is pinned in \texttt{lake-manifest.json}.
\end{thebibliography}
\end{document}
